\section{Connection and obstruction}\label{sec:two-examples}

We use the same small carrier for two contrasting reachability
examples. Take \(X=\{0,1\}\) with its single
nonconstant parity probe and \(R=\{0,1,2,3\}\) with uniform mass.
The four marks have rate \(1/4\) and each reverses to itself.
Take \(\delta=1/8\), \(\kappa=1/2\), horizon two and two grades.
Let \(F_*\) denote the split \(\{0,1\}\mid\{2,3\}\).

The reset register has all entries \(1/4\). The gated register differs
only for one designated mark and the base pair \(g=h=0\). In that
conditional table replace its \(\{0,1\}\)-by-\(\{0,1\}\) submatrix by
\[
 \begin{pmatrix} 1/2 & 0\\ 0&1/2\end{pmatrix}.
\]
Both tables have the required stationary margins and reciprocity.
The conditional mass of each block of \(F_*\) is \(1/2\) for every
preparation in either table. The gated table is genuinely coupled to
the base pair, and loses precisely the two reciprocal proposals
between cells \(0\) and \(1\) at the designated mark and base pair.

\subsection{An exact edit through two zero responses}\label{sec:staged}

For the example of \Cref{sec:opening-example}, use the gated register
at the third mark. Hold the split at both grades, the route equality
between the first--second and first--third words, and the calibration
\(z=1/4\). Fix the fourth response at zero. Every point of
\eqref{eq:opening-path} is a lawful spectrum. Along the first segment
\(y=1/4\); along the last two \(x=0\).
The original route equality and all other tests therefore hold.

For an affine segment with varying response from \(a\) to \(b\),
choose an integer \(D>|b-a|/\varepsilon\) with \(D\ge1\) and sample
\(a+(b-a)k/D\) for \(0\le k\le D\). The difference in each marked
entry is bounded by the magnitude of the varying response change,
by the explicit two-state formula and the fixed rate and register.
The resulting rational edits are strictly within \(\varepsilon\).
The projected probability of retaining the parity bit at the
second mark changes from \(5/8\) to \(3/8\), conditional on the mark.
The prepared, marked contributions to these probabilities are
\(5/32\) and \(3/32\), respectively.
Throughout this construction the native support is fixed, since
\(Q_m>0\) and the register has the same support.

The held calibration requires \(w_z=0\) on every support containing
\(z\), and the reversal rows are identities. On
\(\{x,y,z\}\), the held word requires \(w_y=w_z\).
It follows that a single contraction cannot reach the core
\(\{z\}\), and that a surviving \(y\) keeps its positive sign.
The legal successive contractions remove \(x\) and then \(y\).
On the target support \(\{y,z\}\) the same second contraction
removes the negative \(y\). Thus both signed cores have the value
\(+1\) at \(z\) and zero elsewhere.
A rule that keeps the initial weight balances at all later supports
forces \(w_y=0\) at every step. Every reduct of each endpoint would
retain its different sign at \(y\). No common signed reduct is possible
under that frozen rule. This failure occurs even though the original
word probabilities are equal throughout the explicit chain.


\subsection{A positive-tolerance obstruction}\label{sec:obstruction}

Name the four marks \(a,b,c,d\) and take the responses
\((1/4,1/4,-1/8,1/8)\), the reset register and the indiscrete
partition at both grades. Hold the route \(aa\sim bb\) and the
calibration \(\alpha_b=1/4\). Compare this law with the one obtained
by replacing only \(\alpha_a=1/4\) by \(-1/4\). Both endpoint laws
pass the same held tests. Their common calibration pins \(b\) at \(1/4\) and their
route equality forces \(\alpha_a^2=\alpha_b^2\). Thus \(w_b=0\)
and \(2w_a=2w_b\) on the occupied supports, so the respective
opposite signs of \(a\) survive. In particular the signed cores
differ. The obstruction is also quantitative. Each intermediate
law must pass the original route and calibration, hence its
\(a\)-response belongs to \(\{1/4,-1/4\}\). With \(N=8\) and
\(\delta=1/8\), a rewrite at tolerance
\[
 \frac{\delta(1/2)}{2N+1}=\frac1{272}
\]
has response change less than \(1/2\) by \eqref{eq:inverse-modulus}.
A discrete step fixes the response. Therefore no chain at this
positive tolerance can cross the two branches.
This contrasts with the earlier example in which a zero
response changes the active damping system without invalidating
the held tests.
